\chapter*{Abstract}
\addcontentsline{toc}{chapter}{Abstract}

Populist parties often campaign against corrupt elites, yet governments that concentrate political
power may also weaken oversight and facilitate institutional capture. We examine this tension using
the V-Dem Country-Year v16 and V-Party V2 datasets. Party-level expert assessments are aggregated to
the senior governing party, filtered for coder reliability, expanded between elections, and joined to
annual country data. The resulting unbalanced panel contains 3,965 country-years for 96 countries
between 1970 and 2019.

The analysis separates three questions that are frequently conflated: contemporaneous association,
temporal ordering, and institutional heterogeneity. First, two-way fixed-effects models compare each
country with itself while absorbing common year shocks and controlling for economic development and
electoral democracy. Second, bidirectional lead-lag and Granger-style specifications, placebo leads,
event studies, first differences, and election-based samples probe whether populism precedes
corruption or vice versa. Third, the composite corruption index is decomposed into executive,
public-sector, legislative, and judicial dimensions, with false-discovery-rate correction.

The positive pooled association between populist government and corruption is imprecise after panel
dependence is considered and reverses sign within countries. The controlled same-year estimate is
small and negative, statistically significant in the main specification but not under every
robustness design. We find no robust temporal direction for the composite index. Decomposition shows
that the negative contemporaneous association is strongest for judicial corruption. The only positive
temporal signal is a small two-year populism-to-legislative-corruption coefficient; it survives the
primary Granger and multiplicity tests but not first-difference or election-only checks. Overall, the
data do not support a broad claim that populist government systematically raises political corruption.
They instead reveal a nuanced, institution-specific, and fragile pattern that warrants cautious
interpretation rather than a causal conclusion.

\textbf{Keywords:} populism; political corruption; panel data; fixed effects; Granger-style tests;
V-Dem; V-Party; institutional capture.
\clearpage
